% !TEX program = xelatex

\documentclass[11pt,a4paper]{ctexart}

% ====== 字体 ======
\usepackage{fontspec}
\setmainfont{Latin Modern Roman}
\usepackage{unicode-math}
% 修改：XITS Math 拥有原生粗体数学字形，解决 \symbfit 加粗无效的问题
\setmathfont{XITS Math}

% ====== 颜色支持 ======
\usepackage{xcolor}

% ====== 页面设置 ======
\usepackage[margin=2.5cm]{geometry}
\usepackage{setspace}
\onehalfspacing
\pagestyle{plain}

% ====== 数学宏包 ======
\usepackage{mathtools}
\usepackage{amsmath}
% \usepackage{amssymb}   % 删除
% \usepackage{bm}        % 彻底删除

% ====== 三线表 + 数字对齐 ======
\usepackage{booktabs}
\usepackage{siunitx}
\usepackage{enumitem}

% ====== 交叉引用与超链接 ======
\usepackage[colorlinks=true,
            linkcolor=blue,
            citecolor=green,
            urlcolor=cyan]{hyperref}
\usepackage[noabbrev,nameinlink]{cleveref}

% ====== 框架图形 ======
\usepackage{tcolorbox}  % 引入 tcolorbox 宏包
\tcbuselibrary{skins, breakable} % 加载所需库，以实现更多样式和自动分页

% 定义一个名为 block 的环境
\newtcolorbox{block}{
    colback=white,        % 背景色
    colframe=black!75,    % 边框颜色
    boxrule=0.5mm,        % 边框粗细
    arc=2mm,              % 圆角大小
    boxsep=5pt,           % 边框与内容的间距
    left=6pt, right=6pt,  % 左右内边距
    top=6pt, bottom=6pt,  % 上下内边距
    breakable             % 允许盒子在跨页时自动断开
}

% ====== 自定义命令 ======
\newcommand{\dd}{\mathrm{d}}
\newcommand{\source}[1]{\par \hfill ——#1}

\title{周期性电荷分布的场}
\author{王晨暄}
\date{\today}

\begin{document}
\maketitle

在本文中我们来考虑两种有趣的周期性的电荷分布体系，并且针对这两种情况具体计算空间中的非奇点区域的电势分布。
除了文中给出的分离变量法的解法之外，还可以考虑积分变换在这一类问题中的应用。

\section{无限长带电直线阵列}

\begin{block}
    {\kaishu 考虑无限长的带电直线阵列，电荷分布可以表示为：

    \begin{equation}
        \rho(\symbfit{r})=\sum_{n=-\infty}^{\infty}\lambda\delta(x-na)\delta(z)
        \label{eq:charge_distribution1}
    \end{equation}

    求解空间中的电势函数，表示成级数形式。
    }

\end{block}

\bigskip

首先分析问题：上述电荷分布描述的是平行于 y 轴的无限长带电直线阵列，沿 x 轴方向周期性分布。
所以形成的电场在 x 轴方向上具有周期性，在 y 轴上具有平移不变性。同时，由于电荷分布关于$z=0$以及$x=0$对称，
则有: $\phi(x,z)=\phi(-x,z),\phi(x,-z)=\phi(x,z),\phi(x,z)=\phi(x+a,z)$。

\bigskip

在$z\neq 0$平面上，电势函数满足二维的 Laplace 方程：

\begin{equation}
    \frac{\partial^2\phi}{\partial x^2}+\frac{\partial^2\phi}{\partial z^2}=0
    \label{eq:Laplace}
\end{equation}

分离变量法得到方程的级数形式的解为：

\begin{equation}
    \phi(x,z)=A+B\left|z\right|+\sum_{\gamma}C_{\gamma}\cos(\gamma x)e^{-\gamma\left|z\right|}
    \label{eq:solution}
\end{equation}

考虑到 x 方向的周期性，$\gamma$ 的取值为 $\gamma=\dfrac{2\pi m}{a},m=1,2,3,\cdots$。代入方程\eqref{eq:solution}，得到：

\begin{equation}
    \phi(x,z)=A+B\left|z\right|+\sum_{m=1}^{\infty}C_{m}\cos\left(\frac{2\pi m}{a} x\right)e^{-2\pi m\left|z\right|/a}
\end{equation}

计算$z=0$平面上的法向电场的变化得到：

\begin{align}
    E_z(x,0^+)-E_z(x,0^-)&=-\frac{\partial\phi}{\partial z}\Big|_{z=0^+}+\frac{\partial\phi}{\partial z}\Big|_{z=0^-} \notag\\
    &=-2B+\sum_{m=1}^{\infty}C_{m}\frac{4\pi m}{a}\cos\left(\frac{2\pi m}{a} x\right) \notag
    \label{eq:Ez_jump}
\end{align}

根据边界条件得到：$E_z(x,0^+)-E_z(x,0^-)=\dfrac{\sigma(z,0)}{\epsilon_0}$，带入得到方程：

\begin{equation}
    -2B+\sum_{m=1}^{\infty}C_{m}\frac{4\pi m}{a}\cos\left(\frac{2\pi m}{a} x\right)=\frac{\lambda}{\epsilon_0}\sum_{n=-\infty}^{\infty}\delta(x-na)
    \label{eq:boundary_condition}
\end{equation}

考虑到周期为$a$，将右侧函数在区间$(-\dfrac{a}{2},\dfrac{a}{2})$上展开为 Fourier 级数, 与左边对比系数得到：

\begin{align}
    B&=-\frac{\lambda}{2a\epsilon_0} \notag\\
    C_m&=\frac{\lambda}{2\pi\epsilon_0}\frac{1}{m} \notag
\end{align}

代入方程\eqref{eq:solution}，得到最终的电势函数为：

\begin{equation}
    \phi(x,z)=-\frac{\lambda}{2a\epsilon_0}\left|z\right|+\frac{\lambda}{2\pi\epsilon_0}\sum_{m=1}^{\infty}\frac{1}{m}\cos\left(\frac{2\pi m}{a} x\right)e^{-2\pi m\left|z\right|/a}
    \label{eq:final_solution}
\end{equation}

\bigskip

注意到在$\dfrac{\lambda}{a}$ 有限，$a \to 0$时，电荷分布趋向于面电荷密度为$\sigma=\dfrac{\lambda}{a}$的无限大平面，
此时方程\eqref{eq:final_solution}中的级数项趋向于零，
电势函数趋向于$\phi(x,z)=-\dfrac{\sigma}{2\epsilon_0}\left|z\right|$，与无限大平面电荷分布的结果一致。

\end{document}